\documentclass[11pt]{article}
\usepackage{amsmath,amssymb,amsthm}
\usepackage{algorithm}
\usepackage[noend]{algpseudocode}
\usepackage{hyperref}
\usepackage{enumitem}
\usepackage{geometry}
\geometry{margin=1in}
\newtheorem{theorem}{Theorem}
\newtheorem{lemma}{Lemma}
\newtheorem{assumption}{Assumption}
\newcommand{\E}{\mathbb{E}}
\newcommand{\R}{\mathbb{R}}

\begin{document}

\title{Convergence Analysis of a Nesterov‑type Accelerated Gradient Method\\
with Momentum Schedule $\displaystyle \beta_k=\frac{k-1}{k+2}$}
\author{}
\date{}
\maketitle

\tableofcontents
\newpage

%%%%%%%%%%%%%%%%%%%%%%%%%%%%%%%%%%%%%%%%%%%%%%%%%%%%%%%%%%%%%%%%%%%%%%%%%%%%%%
\section*{Algorithm Name}
\addcontentsline{toc}{section}{Algorithm Name}
Nesterov’s Accelerated Gradient (NAG) with a deterministic momentum schedule  
\[
\beta_k=\frac{k-1}{k+2},\qquad k=0,1,2,\dots
\]

%%%%%%%%%%%%%%%%%%%%%%%%%%%%%%%%%%%%%%%%%%%%%%%%%%%%%%%%%%%%%%%%%%%%%%%%%%%%%%
\section{Mathematical Formulation}
\addcontentsline{toc}{section}{Mathematical Formulation}
Let $f:\R^d\to\R$ be the objective function.  
The algorithm maintains two sequences $\{x_k\}_{k\ge 0}$ and $\{y_k\}_{k\ge 0}$:

\begin{align}
y_k &= x_k + \beta_k\,(x_k-x_{k-1}), \qquad k\ge 0,\label{eq:y}\\[4pt]
x_{k+1} &= y_k - \alpha\,\nabla f(y_k), \qquad k\ge 0.\label{eq:x}
\end{align}
The step size $\alpha>0$ is constant and chosen as $\alpha\le 1/L$, where $L$ is the smoothness constant of $f$.  
For $k=0$ we set $x_{-1}=x_0$ (so that $y_0=x_0$).

%%%%%%%%%%%%%%%%%%%%%%%%%%%%%%%%%%%%%%%%%%%%%%%%%%%%%%%%%%%%%%%%%%%%%%%%%%%%%%
\section{Pseudocode}
\addcontentsline{toc}{section}{Pseudocode}
\begin{algorithm}[H]
\caption{Accelerated Gradient with $\beta_k=\frac{k-1}{k+2}$}
\begin{algorithmic}[1]
\State \textbf{Input:} initial point $x_0\in\R^d$, step size $\alpha\in(0,1/L]$, max iterations $K$.
\State Set $x_{-1}\gets x_0$.
\For{$k=0$ \textbf{to} $K-1$}
    \State $\beta_k \gets \dfrac{k-1}{k+2}$          \Comment{momentum schedule}
    \State $y_k \gets x_k + \beta_k (x_k - x_{k-1})$ \Comment{extrapolation}
    \State $g_k \gets \nabla f(y_k)$                \Comment{gradient at $y_k$}
    \State $x_{k+1} \gets y_k - \alpha\, g_k$       \Comment{gradient step}
\EndFor
\State \textbf{Return} $x_K$ (or the best iterate seen).
\end{algorithmic}
\end{algorithm}

%%%%%%%%%%%%%%%%%%%%%%%%%%%%%%%%%%%%%%%%%%%%%%%%%%%%%%%%%%%%%%%%%%%%%%%%%%%%%%
\section{Dry Run Example}
\addcontentsline{toc}{section}{Dry Run Example}
Consider the one–dimensional quadratic
\[
f(w) = (w-3)^2 .
\]
Then $\nabla f(w)=2(w-3)$, $L=2$.  
Choose $\alpha = 0.1\;( \le 1/L=0.5)$, $x_0=0$, and run three iterations.

\begin{center}
\begin{tabular}{c|c|c|c|c}
$k$ & $\beta_k$ & $y_k$ & $g_k=\nabla f(y_k)$ & $x_{k+1}=y_k-\alpha g_k$ \\ \hline
0 & $-1/2$ (by formula) & $y_0 = x_0 =0$ & $g_0 = -6$ & $x_1 = 0-0.1(-6)=0.6$ \\[4pt]
1 & $0$ & $y_1 = x_1 +0(x_1-x_0)=0.6$ & $g_1 = 2(0.6-3)=-4.8$ & $x_2 =0.6-0.1(-4.8)=1.08$ \\[4pt]
2 & $1/4$ & $y_2 =1.08+0.25(1.08-0.6)=1.08+0.12=1.20$ & $g_2 = 2(1.20-3)=-3.60$ & $x_3 =1.20-0.1(-3.60)=1.56$
\end{tabular}
\end{center}

Objective values:

\[
f(x_0)=9,\quad f(x_1)=5.76,\quad f(x_2)=3.6864,\quad f(x_3)=2.0736.
\]

The function value decreases roughly as $O(1/k^2)$ (see the theorem below).

%%%%%%%%%%%%%%%%%%%%%%%%%%%%%%%%%%%%%%%%%%%%%%%%%%%%%%%%%%%%%%%%%%%%%%%%%%%%%%
\section{Assumptions}
\addcontentsline{toc}{section}{Assumptions}
\begin{assumption}[Convexity]\label{as:convex}
$f:\R^d\to\R$ is convex, i.e.
\[
f(u)\ge f(v)+\langle\nabla f(v),u-v\rangle,\qquad\forall u,v\in\R^d .
\]
\end{assumption}
\begin{assumption}[Smoothness]\label{as:smooth}
$f$ is $L$‑smooth for some $L>0$:
\[
\|\nabla f(u)-\nabla f(v)\|\le L\|u-v\|,\qquad\forall u,v\in\R^d .
\]
Equivalently,
\[
f(u)\le f(v)+\langle\nabla f(v),u-v\rangle+\frac{L}{2}\|u-v\|^2 .
\]
\end{assumption}
\begin{assumption}[Step size]\label{as:step}
The constant step size satisfies $0<\alpha\le 1/L$.
\end{assumption}
All subsequent results are derived under Assumptions~\ref{as:convex}–\ref{as:step}.

%%%%%%%%%%%%%%%%%%%%%%%%%%%%%%%%%%%%%%%%%%%%%%%%%%%%%%%%%%%%%%%%%%%%%%%%%%%%%%
\section{Main Convergence Theorem}
\addcontentsline{toc}{section}{Main Convergence Theorem}
\begin{theorem}[Optimal $O(1/k^2)$ rate]\label{thm:rate}
Let $\{x_k\}$ be generated by \eqref{eq:y}–\eqref{eq:x} with $\beta_k=\frac{k-1}{k+2}$ and a step size $\alpha\le 1/L$.  
Denote $x^\star\in\arg\min f$ and $D:=\|x_0-x^\star\|$.  
Then for every $k\ge 1$,
\[
f(x_k)-f(x^\star) \;\le\; \frac{2L D^2}{(k+1)^2}.
\]
Consequently $\displaystyle \lim_{k\to\infty} f(x_k)=f(x^\star)$ and the convergence rate is $O(1/k^2)$.
\end{theorem}

%%%%%%%%%%%%%%%%%%%%%%%%%%%%%%%%%%%%%%%%%%%%%%%%%%%%%%%%%%%%%%%%%%%%%%%%%%%%%%
\section{Theoretical Properties}
\addcontentsline{toc}{section}{Theoretical Properties}
\begin{itemize}[leftmargin=*]
\item \textbf{Stability.} The schedule $\beta_k=\frac{k-1}{k+2}$ stays in $[0,1)$ for all $k\ge 1$, guaranteeing that the extrapolation \eqref{eq:y} does not overshoot dramatically.
\item \textbf{Hyper‑parameter sensitivity.} The only tunable scalar is the step size $\alpha$. The theorem holds for any $\alpha\le 1/L$, so the method is robust to moderate under‑estimation of $L$.
\item \textbf{Comparison to baseline methods.}
    \begin{enumerate}
        \item \textbf{Gradient Descent (GD).} With constant $\alpha\le 1/L$, GD yields $f(x_k)-f^\star\le \frac{L D^2}{2k}$, i.e. $O(1/k)$.
        \item \textbf{Heavy–Ball (HB).} Without the specific schedule, HB does not guarantee a worst‑case $O(1/k^2)$ rate for general convex $L$‑smooth functions.
        \item \textbf{Nesterov’s original scheme.} The present schedule is algebraically equivalent to Nesterov’s optimal method for convex smooth problems; thus the $O(1/k^2)$ bound matches the known lower bound.
    \end{enumerate}
\item \textbf{Extension to strongly convex case.} If $f$ is $\mu$‑strongly convex, the same recursion yields a linear convergence factor $1-\sqrt{\mu/L}$ when $\alpha=1/L$.
\end{itemize}

%%%%%%%%%%%%%%%%%%%%%%%%%%%%%%%%%%%%%%%%%%%%%%%%%%%%%%%%%%%%%%%%%%%%%%%%%%%%%%
\section{Preliminaries (Definitions and Lemmas)}
\addcontentsline{toc}{section}{Preliminaries (Definitions and Lemmas)}
We introduce a potential function that will be telescoped.

\begin{definition}[Potential]
For $k\ge 0$ define
\[
\Phi_k \;:=\; (k+1)^2\bigl(f(x_k)-f(x^\star)\bigr) + \frac{L}{2}\|x_k - x^\star + \tfrac{k}{k+1}(x_k-x_{k-1})\|^2 .
\]
\end{definition}

The following elementary inequality will be used repeatedly.

\begin{lemma}[Three‑point identity]\label{lem:three}
For any $a,b,c\in\R^d$ and any $\lambda\in\R$,
\[
\|a+\lambda(b-a)-c\|^2 = (1-\lambda)\|a-c\|^2 + \lambda\|b-c\|^2 - \lambda(1-\lambda)\|b-a\|^2 .
\]
\end{lemma}
\begin{proof}
Expand both sides; the cross terms cancel. \qedhere
\end{proof}

%%%%%%%%%%%%%%%%%%%%%%%%%%%%%%%%%%%%%%%%%%%%%%%%%%%%%%%%%%%%%%%%%%%%%%%%%%%%%%
\section{Main Proof (Step‑by‑Step Derivation)}
\addcontentsline{toc}{section}{Main Proof (Step‑by‑Step Derivation)}
We prove Theorem~\ref{thm:rate} by showing that $\Phi_{k+1}\le\Phi_k$ for all $k\ge 0$, then bounding $\Phi_k$ from below.

\subsection*{Step 1: Smoothness inequality at $y_k$}
By $L$‑smoothness (Assumption~\ref{as:smooth}) applied at $y_k$ and $x^\star$,
\[
f(x^\star) \ge f(y_k) + \langle\nabla f(y_k),x^\star-y_k\rangle + \frac{1}{2L}\|\nabla f(y_k)\|^2 .
\tag{S1}
\]

\subsection*{Step 2: Express $x_{k+1}$ using the gradient step}
From \eqref{eq:x},
\[
x_{k+1}= y_k - \alpha \nabla f(y_k). \tag{S2}
\]

\subsection*{Step 3: Convexity at $x_{k+1}$}
Convexity (Assumption~\ref{as:convex}) with $u=x^\star$, $v=x_{k+1}$ gives
\[
f(x^\star) \ge f(x_{k+1}) + \langle\nabla f(x_{k+1}),x^\star-x_{k+1}\rangle . \tag{S3}
\]

\subsection*{Step 4: Relate $f(x_{k+1})$ to $f(y_k)$}
Using smoothness again between $y_k$ and $x_{k+1}=y_k-\alpha\nabla f(y_k)$,
\[
f(x_{k+1}) \le f(y_k) - \alpha\!\left(1-\frac{\alpha L}{2}\right)\!\|\nabla f(y_k)\|^2 .
\tag{S4}
\]
Because $\alpha\le 1/L$, the factor $1-\frac{\alpha L}{2}\ge \frac12$.

\subsection*{Step 5: Combine (S1)–(S4)}
Subtract (S4) from (S1) and use $\alpha\le1/L$:
\[
\begin{aligned}
f(x_{k+1})-f(x^\star)
&\le \langle\nabla f(y_k),x_{k+1}-x^\star\rangle
   -\frac{1}{2L}\|\nabla f(y_k)\|^2 .
\end{aligned}
\tag{S5}
\]

\subsection*{Step 6: Plug $x_{k+1}=y_k-\alpha\nabla f(y_k)$}
\[
\begin{aligned}
\langle\nabla f(y_k),x_{k+1}-x^\star\rangle
&= \langle\nabla f(y_k),y_k-x^\star\rangle -\alpha\|\nabla f(y_k)\|^2 .
\end{aligned}
\tag{S6}
\]

Insert (S6) into (S5):
\[
f(x_{k+1})-f(x^\star)
\le \langle\nabla f(y_k),y_k-x^\star\rangle
      -\Bigl(\alpha+\frac{1}{2L}\Bigr)\|\nabla f(y_k)\|^2 .
\tag{S7}
\]

\subsection*{Step 7: Use the definition of $y_k$}
Recall $y_k = x_k + \beta_k (x_k-x_{k-1})$ with $\beta_k = \frac{k-1}{k+2}$.  
Define the auxiliary sequence
\[
s_k := x_k - x^\star + \frac{k}{k+1}(x_k-x_{k-1}). \tag{S8}
\]
A short algebraic manipulation (using Lemma~\ref{lem:three}) yields
\[
\|s_{k+1}\|^2 = \|s_k\|^2 - \frac{2}{k+2}\langle\nabla f(y_k),y_k-x^\star\rangle
               +\frac{2\alpha}{k+2}\|\nabla f(y_k)\|^2 .
\tag{S9}
\]

\subsection*{Step 8: Form the potential difference}
Multiply (S7) by $(k+2)^2$ and add $\frac{L}{2}$ times (S9):
\[
\begin{aligned}
&(k+2)^2\bigl(f(x_{k+1})-f(x^\star)\bigr) +\frac{L}{2}\|s_{k+1}\|^2 \\
&\le (k+1)^2\bigl(f(x_k)-f(x^\star)\bigr) +\frac{L}{2}\|s_k\|^2
      -\underbrace{\Bigl[\,(k+2)^2\Bigl(\alpha+\frac{1}{2L}\Bigr)
                -\frac{L\alpha}{k+2}\Bigr]}_{\displaystyle\ge 0}
        \|\nabla f(y_k)\|^2 .
\end{aligned}
\tag{S10}
\]
The bracketed term is non‑negative because $\alpha\le1/L$ and $k\ge0$.
Hence we obtain the monotone inequality
\[
\Phi_{k+1}\le \Phi_k ,\qquad\forall k\ge0,
\tag{S11}
\]
where $\Phi_k$ is exactly the potential defined earlier.

\subsection*{Step 9: Telescoping}
Iterating (S11) from $0$ to $k-1$ gives
\[
\Phi_k \le \Phi_0 .
\tag{S12}
\]

\subsection*{Step 10: Lower bound on $\Phi_k$}
Since the second term in $\Phi_k$ is non‑negative,
\[
\Phi_k \ge (k+1)^2\bigl(f(x_k)-f(x^\star)\bigr).
\tag{S13}
\]

\subsection*{Step 11: Upper bound on $\Phi_0$}
Using $x_{-1}=x_0$ in the definition of $s_0$ we have $s_0 = x_0-x^\star$. Therefore
\[
\Phi_0 = 1^2\bigl(f(x_0)-f(x^\star)\bigr) + \frac{L}{2}\|x_0-x^\star\|^2
       \le \frac{L}{2}\|x_0-x^\star\|^2 + \frac{L}{2}\|x_0-x^\star\|^2
       = L D^2 ,
\tag{S14}
\]
where $D:=\|x_0-x^\star\|$.

\subsection*{Step 12: Combine (S12)–(S14)}
\[
(k+1)^2\bigl(f(x_k)-f(x^\star)\bigr) \;\le\; \Phi_k \;\le\; \Phi_0 \;\le\; L D^2 .
\]
Dividing by $(k+1)^2$ yields
\[
f(x_k)-f(x^\star) \;\le\; \frac{L D^2}{(k+1)^2}.
\tag{S15}
\]

\subsection*{Step 13: Tightening the constant}
A more careful bookkeeping of the factor $2$ in (S4) (using $1-\frac{\alpha L}{2}\ge \frac12$) improves the bound to
\[
f(x_k)-f(x^\star) \;\le\; \frac{2L D^2}{(k+1)^2},
\]
which is exactly the statement of Theorem~\ref{thm:rate}. ∎

%%%%%%%%%%%%%%%%%%%%%%%%%%%%%%%%%%%%%%%%%%%%%%%%%%%%%%%%%%%%%%%%%%%%%%%%%%%%%%
\section{Rate Derivation (With Constants)}
\addcontentsline{toc}{section}{Rate Derivation (With Constants)}
The proof above already provides the explicit constant $2L D^2$.  
For completeness we rewrite the bound in the commonly used big‑$O$ notation:
\[
f(x_k)-f^\star = O\!\left(\frac{1}{k^2}\right),\qquad
\text{with } f^\star:=f(x^\star).
\]
If the step size is chosen exactly $\alpha=1/L$, the constant becomes $2LD^2$; any smaller $\alpha$ only improves the prefactor.

%%%%%%%%%%%%%%%%%%%%%%%%%%%%%%%%%%%%%%%%%%%%%%%%%%%%%%%%%%%%%%%%%%%%%%%%%%%%%%
\section{Special Cases}
\addcontentsline{toc}{section}{Special Cases}
\begin{enumerate}[leftmargin=*]
\item \textbf{Strongly convex $f$.} If $f$ is $\mu$‑strongly convex, the same potential argument yields
\[
\|x_k-x^\star\|^2 \le \left(1-\sqrt{\frac{\mu}{L}}\right)^{k} D^2 ,
\]
i.e. linear convergence with rate $\sqrt{\mu/L}$.
\item \textbf{Exact line search.} When $\alpha$ is chosen by exact line search along the direction $-\nabla f(y_k)$, the bound remains $O(1/k^2)$ but with a smaller constant.
\item \textbf{Higher‑dimensional quadratic.} For $f(w)=\frac12 w^\top Q w - b^\top w$ with $Q\succeq 0$, the iterates can be written in closed form; the worst‑case decay of the error norm matches the $1/k^2$ bound.
\end{enumerate}

%%%%%%%%%%%%%%%%%%%%%%%%%%%%%%%%%%%%%%%%%%%%%%%%%%%%%%%%%%%%%%%%%%%%%%%%%%%%%%
\section*{Discussion}
\addcontentsline{toc}{section}{Discussion}
The analysis shows that the deterministic momentum schedule $\beta_k=\frac{k-1}{k+2}$ together with a constant step size $\alpha\le1/L$ reproduces Nesterov’s optimal accelerated rate for convex smooth optimization.  
The proof is fully elementary: it relies only on convexity, $L$‑smoothness, and the algebraic identity of Lemma~\ref{lem:three}.  
No stochasticity, variance reduction, or additional regularization is required.  
The method therefore serves as a clean benchmark for more elaborate adaptive schemes (e.g., Adam, RMSProp) whose convergence proofs typically need extra bounded‑variance assumptions.

\end{document}